\documentclass[../../main.tex]{subfiles}
\begin{document}

{\bf [解答]}

\begin{figure}[htb]
  \centering
  \begin{tikzpicture}[x={(1.4cm, -0.4cm)}, y={(0.8cm, 0.6cm)}, z={(0cm, 2.2cm)},
    >=stealth, every node/.style={font=\small}]
    % 座標軸
    \draw[->, thick, gray!80!black] (0, 0, 0) -- (1.8, 0, 0) node[below left] {$x$};
    \draw[->, thick, gray!80!black] (0, 0, 0) -- (0, 1.8, 0) node[above right] {$y$};
    \draw[->, thick, gray!80!black] (0, 0, 0) -- (0, 0, 1.6) node[above] {$z$};
    \node[below left=2pt] at (0, 0, 0) {$O$};

    % 頂点
    \coordinate (V) at (0, 0, 1);
    \coordinate (A) at (1, 1, 0);
    \coordinate (B) at (1, -1, 0);
    \coordinate (C) at (-1, -1, 0);
    \coordinate (D) at (-1, 1, 0);

    % 斜面のハッチング
    \fill[pattern=north east lines, pattern color=blue!40] (V) -- (A) -- (B) -- cycle;
    \fill[pattern=north west lines, pattern color=orange!40] (V) -- (B) -- (C) -- cycle;

    % 奥の稜線（破線）
    \draw[thick, dashed, gray!70!black] (A) -- (D) -- (C);
    \draw[thick, dashed, gray!70!black] (V) -- (D);

    % 手前の稜線（実線）
    \draw[line width=1.1pt, gray!85!black] (C) -- (B) -- (A);
    \draw[line width=1.2pt, blue!85!black] (V) -- (A);
    \draw[line width=1.3pt, purple!85!black] (V) -- (B);
    \draw[line width=1.2pt, orange!85!black] (V) -- (C);

    % 頂点ノード
    \node[circle, fill=blue!85!black, inner sep=1.2pt] at (V) {};
    \node[above=2pt] at (V) {$V(0,0,1)$};
    \node[circle, fill=blue!85!black, inner sep=1.2pt] at (A) {};
    \node[right=2pt] at (A) {$A(1,1,0)$};
    \node[circle, fill=blue!85!black, inner sep=1.2pt] at (B) {};
    \node[below=2pt] at (B) {$B(1,-1,0)$};
    \node[circle, fill=blue!85!black, inner sep=1.2pt] at (C) {};
    \node[left=2pt] at (C) {$C(-1,-1,0)$};
    \node[circle, fill=gray!70!black, inner sep=1.2pt] at (D) {};
    \node[above left=1pt, gray!70!black] at (D) {$D(-1,1,0)$};
  \end{tikzpicture}
  \caption{正四角錐 $V-ABCD$ と斜面 $VAB$, $VBC$}
  \label{fig:pyramid_v_abcd}
\end{figure}

底面が一辺 $2$ の正方形であるとしてよい．そこで $xyz$ 座標を設定し $A(1,1,0)$，
$B(1,-1,0)$，$C(-1,-1,0)$，$D(-1,1,0)$ とおく．底面と斜面の成す角が $\pi/4$ だから，
$V(0,0,1)$ として考えてよい．対称性から平面 $VAB$ と平面 $VBC$ のみ考えればよい．これらの平面の基底をなす $2$ ベクトルは
\begin{align}
  VAB:&
  \begin{pmatrix}0\\1\\0
  \end{pmatrix} ,
  \begin{pmatrix}-1\\0\\1
  \end{pmatrix}  \\
  VBC:&
  \begin{pmatrix}1\\0\\0
  \end{pmatrix} ,
  \begin{pmatrix}0\\1\\1
  \end{pmatrix}
\end{align}
であるから，法線ベクトルは
\begin{align}
  &VAB:\vec{a}=
  \begin{pmatrix}1\\0\\1
  \end{pmatrix} &VBC:\vec{b}=
  \begin{pmatrix}0\\-1\\1
  \end{pmatrix}
\end{align}
となる．故にこれら $2$ ベクトルの成す角 $\theta$ として ($0\le\theta\le\pi$)
\begin{align}
  \cos\theta&=\dfrac{\vec{a}\cdot\vec{b}}{|\vec{a}||\vec{b}|} \\
  &=\dfrac{1}{\sqrt{2}\sqrt{2}}=\dfrac{1}{2}
\end{align}
である．従って $\theta=\dfrac{\pi}{3}$ である．
故に $2$ 平面の成す角 $t$ は $0\le t\le\dfrac{\pi}{2}$ で考えて
\begin{align}
  t=\min(\theta,\pi-\theta)=\dfrac{\pi}{3}
\end{align}
である．$\cdots$(答)

\end{document}
